\chapter*{Annexes Techniques}
\addcontentsline{toc}{chapter}{Annexes Techniques}

\section*{Annexe A : Dictionnaire des Données et Configuration}
\markboth{ANNEXES TECHNIQUES}{ANNEXE A : DICTIONNAIRE DES DONNÉES}

\subsection*{Configuration de l'Établissement (\texttt{configs/hospital.json})}
L'extrait suivant est produit depuis le fichier de configuration :
\lstinputlisting[language=json, caption={Extrait de la configuration de l'établissement.}]{tables/hospital_extrait.json}

\subsection*{Structure de la Table des Admissions}
Le tableau~\ref{tab:dictionnaire} décrit les colonnes de \texttt{data/synthetic/historical\_demand.csv}.

% Généré par rapport/scripts/generate_tables.py : ne pas modifier à la main.
\begin{table}[H]
    \centering
    \caption{Colonnes de \texttt{historical\_demand.csv}.}
    \label{tab:dictionnaire}
    \small
    \begin{tabular}{lll}
        \toprule
        \textbf{Colonne} & \textbf{Type} & \textbf{Description} \\
        \midrule
        \texttt{date} & Date ISO & Jour observé \\
        \texttt{service} & Texte & Code du service : URG, MED, REA ou CHI \\
        \texttt{admissions} & Entier & Admissions du jour \\
        \texttt{discharges} & Entier & Sorties issues de la dynamique des séjours \\
        \texttt{emergency\_arrivals} & Entier ou vide & Arrivées aux urgences, renseignées pour URG seulement \\
        \texttt{is\_holiday} & Booléen & Jour férié synthétique \\
        \texttt{is\_special\_event} & Booléen & Événement exceptionnel tiré au hasard \\
        \texttt{data\_source} & Texte & Toujours synthetic \\
        \bottomrule
    \end{tabular}
\end{table}


\newpage
\section*{Annexe B : Guide d'Exploitation et Commandes}
\markboth{ANNEXES TECHNIQUES}{ANNEXE B : GUIDE D'EXPLOITATION}

Les commandes sont données pour Windows PowerShell, depuis la racine du projet.

\subsection*{Installation}
\begin{lstlisting}[caption={Installation de l'environnement.}]
python -m venv .venv
.\.venv\Scripts\python.exe -m pip install . pytest pytest-cov ruff httpx
$env:PYTHONPATH = "src"
cd frontend
npm.cmd install
\end{lstlisting}

\subsection*{Chaîne de Calcul}
\begin{lstlisting}[caption={Commandes de la chaîne de calcul.}]
# 1. Donnees synthetiques et rapport de qualite
.\.venv\Scripts\python.exe -m hospital_planning.data.generator

# 2. Comparaison des modeles et previsions J+1 a J+7
.\.venv\Scripts\python.exe -m hospital_planning.forecasting.evaluation

# 3. Instance operationnelle S0
.\.venv\Scripts\python.exe -m hospital_planning.operational.requirements

# 4. Comparaison complete (environ 32 minutes)
.\.venv\Scripts\python.exe -m hospital_planning.experiments.runner --methods naive greedy milp cp_sat --profiles equilibre priorite_accueil priorite_personnel

# 5. Grille de sensibilite, Naive et Greedy (environ 100 secondes)
.\.venv\Scripts\python.exe -m hospital_planning.experiments.sensitivity

# 6. Tableaux et chiffres de ce rapport
.\.venv\Scripts\python.exe rapport/scripts/generate_tables.py
\end{lstlisting}

\subsection*{API et Tableau de Bord}
\begin{lstlisting}[caption={Lancement de l'API et du tableau de bord (deux terminaux).}]
# Terminal 1 : API, documentation sur http://127.0.0.1:8000/docs
.\.venv\Scripts\python.exe -m uvicorn hospital_planning.api.app:app --host 127.0.0.1 --port 8000

# Terminal 2 : tableau de bord sur http://localhost:3011
cd frontend
npm.cmd run dev
\end{lstlisting}

\subsection*{Contrôles}
\begin{lstlisting}[caption={Tests et analyse statique.}]
.\.venv\Scripts\python.exe -m pytest -q
.\.venv\Scripts\python.exe -m ruff check src tests
cd frontend
npm.cmd run typecheck
npm.cmd run lint
npm.cmd run test
\end{lstlisting}

\newpage
\section*{Annexe C : Grille de Sensibilité Complète (80 Cas)}
\markboth{ANNEXES TECHNIQUES}{ANNEXE C : GRILLE DE SENSIBILITÉ}

Le tableau~\ref{tab:sensibilite_complete} donne les résultats de Greedy. Les cas dont les contraintes ne sont pas respectées ne décrivent pas un plan réalisable.

% Généré par rapport/scripts/generate_tables.py : ne pas modifier à la main.
{\footnotesize
\begin{longtable}{lcccrrrrc}
\caption{Les 80 cas de la grille de sensibilité, heuristique Greedy, profil équilibré.}\label{tab:sensibilite_complete}\\
\toprule
\textbf{Cas} & \textbf{Abs.} & \textbf{Dem.} & \textbf{Lits} & \textbf{Objectif} & \textbf{Refus} & \textbf{Déficit RH} & \textbf{Postes couv.} & \textbf{Contraintes} \\
\midrule\endfirsthead
\toprule
\textbf{Cas} & \textbf{Abs.} & \textbf{Dem.} & \textbf{Lits} & \textbf{Objectif} & \textbf{Refus} & \textbf{Déficit RH} & \textbf{Postes couv.} & \textbf{Contraintes} \\
\midrule\endhead
\bottomrule\endfoot
\texttt{A00\_D100\_C00} & 0\,\% & $\times$1{,}00 & -- & 636\,379 & 4 & 1\,200 & 48{,}3\,\% & respectées \\
\texttt{A00\_D100\_C05} & 0\,\% & $\times$1{,}00 & $-$5\,\% & 650\,419 & 14 & 1\,137 & 49{,}7\,\% & respectées \\
\texttt{A00\_D100\_C10} & 0\,\% & $\times$1{,}00 & $-$10\,\% & 674\,469 & 26 & 1\,096 & 50{,}6\,\% & respectées \\
\texttt{A00\_D100\_C20} & 0\,\% & $\times$1{,}00 & $-$20\,\% & 726\,459 & 79 & 901 & 55{,}3\,\% & non \\
\texttt{A00\_D125\_C00} & 0\,\% & $\times$1{,}25 & -- & 924\,327 & 73 & 1\,408 & 44{,}6\,\% & respectées \\
\texttt{A00\_D125\_C05} & 0\,\% & $\times$1{,}25 & $-$5\,\% & 967\,887 & 105 & 1\,309 & 46{,}4\,\% & respectées \\
\texttt{A00\_D125\_C10} & 0\,\% & $\times$1{,}25 & $-$10\,\% & 997\,969 & 139 & 1\,187 & 48{,}6\,\% & respectées \\
\texttt{A00\_D125\_C20} & 0\,\% & $\times$1{,}25 & $-$20\,\% & 1\,079\,619 & 206 & 984 & 53{,}2\,\% & non \\
\texttt{A00\_D150\_C00} & 0\,\% & $\times$1{,}50 & -- & 1\,246\,627 & 188 & 1\,450 & 43{,}9\,\% & respectées \\
\texttt{A00\_D150\_C05} & 0\,\% & $\times$1{,}50 & $-$5\,\% & 1\,284\,547 & 227 & 1\,312 & 46{,}3\,\% & respectées \\
\texttt{A00\_D150\_C10} & 0\,\% & $\times$1{,}50 & $-$10\,\% & 1\,318\,589 & 262 & 1\,190 & 48{,}6\,\% & respectées \\
\texttt{A00\_D150\_C20} & 0\,\% & $\times$1{,}50 & $-$20\,\% & 1\,403\,209 & 329 & 987 & 53{,}1\,\% & non \\
\texttt{A00\_D200\_C00} & 0\,\% & $\times$2{,}00 & -- & 1\,906\,827 & 434 & 1\,468 & 43{,}6\,\% & respectées \\
\texttt{A00\_D200\_C05} & 0\,\% & $\times$2{,}00 & $-$5\,\% & 1\,948\,707 & 473 & 1\,330 & 46{,}1\,\% & respectées \\
\texttt{A00\_D200\_C10} & 0\,\% & $\times$2{,}00 & $-$10\,\% & 1\,984\,219 & 508 & 1\,205 & 48{,}2\,\% & respectées \\
\texttt{A00\_D200\_C20} & 0\,\% & $\times$2{,}00 & $-$20\,\% & 2\,076\,279 & 575 & 1\,005 & 52{,}7\,\% & non \\
\texttt{A05\_D100\_C00} & 5\,\% & $\times$1{,}00 & -- & 660\,879 & 4 & 1\,249 & 46{,}2\,\% & respectées \\
\texttt{A05\_D100\_C05} & 5\,\% & $\times$1{,}00 & $-$5\,\% & 674\,919 & 14 & 1\,186 & 47{,}5\,\% & respectées \\
\texttt{A05\_D100\_C10} & 5\,\% & $\times$1{,}00 & $-$10\,\% & 698\,469 & 26 & 1\,144 & 48{,}4\,\% & respectées \\
\texttt{A05\_D100\_C20} & 5\,\% & $\times$1{,}00 & $-$20\,\% & 748\,959 & 79 & 946 & 53{,}1\,\% & non \\
\texttt{A05\_D125\_C00} & 5\,\% & $\times$1{,}25 & -- & 947\,327 & 73 & 1\,454 & 42{,}8\,\% & respectées \\
\texttt{A05\_D125\_C05} & 5\,\% & $\times$1{,}25 & $-$5\,\% & 990\,887 & 105 & 1\,355 & 44{,}5\,\% & respectées \\
\texttt{A05\_D125\_C10} & 5\,\% & $\times$1{,}25 & $-$10\,\% & 1\,020\,969 & 139 & 1\,233 & 46{,}6\,\% & respectées \\
\texttt{A05\_D125\_C20} & 5\,\% & $\times$1{,}25 & $-$20\,\% & 1\,103\,119 & 206 & 1\,031 & 51{,}0\,\% & non \\
\texttt{A05\_D150\_C00} & 5\,\% & $\times$1{,}50 & -- & 1\,269\,627 & 188 & 1\,496 & 42{,}1\,\% & respectées \\
\texttt{A05\_D150\_C05} & 5\,\% & $\times$1{,}50 & $-$5\,\% & 1\,307\,547 & 227 & 1\,358 & 44{,}5\,\% & respectées \\
\texttt{A05\_D150\_C10} & 5\,\% & $\times$1{,}50 & $-$10\,\% & 1\,341\,589 & 262 & 1\,236 & 46{,}6\,\% & respectées \\
\texttt{A05\_D150\_C20} & 5\,\% & $\times$1{,}50 & $-$20\,\% & 1\,426\,709 & 329 & 1\,034 & 50{,}9\,\% & non \\
\texttt{A05\_D200\_C00} & 5\,\% & $\times$2{,}00 & -- & 1\,929\,327 & 434 & 1\,513 & 41{,}9\,\% & respectées \\
\texttt{A05\_D200\_C05} & 5\,\% & $\times$2{,}00 & $-$5\,\% & 1\,971\,207 & 473 & 1\,375 & 44{,}2\,\% & respectées \\
\texttt{A05\_D200\_C10} & 5\,\% & $\times$2{,}00 & $-$10\,\% & 2\,007\,219 & 508 & 1\,251 & 46{,}3\,\% & respectées \\
\texttt{A05\_D200\_C20} & 5\,\% & $\times$2{,}00 & $-$20\,\% & 2\,099\,779 & 575 & 1\,052 & 50{,}5\,\% & non \\
\texttt{A10\_D100\_C00} & 10\,\% & $\times$1{,}00 & -- & 685\,877 & 4 & 1\,299 & 44{,}1\,\% & respectées \\
\texttt{A10\_D100\_C05} & 10\,\% & $\times$1{,}00 & $-$5\,\% & 699\,917 & 14 & 1\,236 & 45{,}3\,\% & respectées \\
\texttt{A10\_D100\_C10} & 10\,\% & $\times$1{,}00 & $-$10\,\% & 723\,467 & 26 & 1\,194 & 46{,}1\,\% & respectées \\
\texttt{A10\_D100\_C20} & 10\,\% & $\times$1{,}00 & $-$20\,\% & 772\,957 & 79 & 994 & 50{,}7\,\% & non \\
\texttt{A10\_D125\_C00} & 10\,\% & $\times$1{,}25 & -- & 972\,825 & 73 & 1\,505 & 40{,}8\,\% & respectées \\
\texttt{A10\_D125\_C05} & 10\,\% & $\times$1{,}25 & $-$5\,\% & 1\,016\,385 & 105 & 1\,406 & 42{,}4\,\% & respectées \\
\texttt{A10\_D125\_C10} & 10\,\% & $\times$1{,}25 & $-$10\,\% & 1\,046\,967 & 139 & 1\,285 & 44{,}4\,\% & respectées \\
\texttt{A10\_D125\_C20} & 10\,\% & $\times$1{,}25 & $-$20\,\% & 1\,127\,617 & 206 & 1\,080 & 48{,}6\,\% & non \\
\texttt{A10\_D150\_C00} & 10\,\% & $\times$1{,}50 & -- & 1\,295\,125 & 188 & 1\,547 & 40{,}1\,\% & respectées \\
\texttt{A10\_D150\_C05} & 10\,\% & $\times$1{,}50 & $-$5\,\% & 1\,333\,045 & 227 & 1\,409 & 42{,}4\,\% & respectées \\
\texttt{A10\_D150\_C10} & 10\,\% & $\times$1{,}50 & $-$10\,\% & 1\,367\,587 & 262 & 1\,288 & 44{,}3\,\% & respectées \\
\texttt{A10\_D150\_C20} & 10\,\% & $\times$1{,}50 & $-$20\,\% & 1\,451\,207 & 329 & 1\,083 & 48{,}6\,\% & non \\
\texttt{A10\_D200\_C00} & 10\,\% & $\times$2{,}00 & -- & 1\,955\,825 & 434 & 1\,566 & 39{,}9\,\% & respectées \\
\texttt{A10\_D200\_C05} & 10\,\% & $\times$2{,}00 & $-$5\,\% & 1\,997\,705 & 473 & 1\,428 & 42{,}1\,\% & respectées \\
\texttt{A10\_D200\_C10} & 10\,\% & $\times$2{,}00 & $-$10\,\% & 2\,032\,716 & 508 & 1\,302 & 44{,}1\,\% & respectées \\
\texttt{A10\_D200\_C20} & 10\,\% & $\times$2{,}00 & $-$20\,\% & 2\,124\,277 & 575 & 1\,101 & 48{,}2\,\% & non \\
\texttt{A15\_D100\_C00} & 15\,\% & $\times$1{,}00 & -- & 712\,876 & 4 & 1\,353 & 41{,}7\,\% & respectées \\
\texttt{A15\_D100\_C05} & 15\,\% & $\times$1{,}00 & $-$5\,\% & 726\,916 & 14 & 1\,290 & 42{,}9\,\% & respectées \\
\texttt{A15\_D100\_C10} & 15\,\% & $\times$1{,}00 & $-$10\,\% & 750\,466 & 26 & 1\,248 & 43{,}7\,\% & respectées \\
\texttt{A15\_D100\_C20} & 15\,\% & $\times$1{,}00 & $-$20\,\% & 800\,956 & 79 & 1\,050 & 47{,}9\,\% & non \\
\texttt{A15\_D125\_C00} & 15\,\% & $\times$1{,}25 & -- & 1\,002\,325 & 73 & 1\,564 & 38{,}4\,\% & respectées \\
\texttt{A15\_D125\_C05} & 15\,\% & $\times$1{,}25 & $-$5\,\% & 1\,045\,885 & 105 & 1\,465 & 40{,}0\,\% & respectées \\
\texttt{A15\_D125\_C10} & 15\,\% & $\times$1{,}25 & $-$10\,\% & 1\,073\,965 & 139 & 1\,339 & 42{,}0\,\% & respectées \\
\texttt{A15\_D125\_C20} & 15\,\% & $\times$1{,}25 & $-$20\,\% & 1\,155\,616 & 206 & 1\,136 & 46{,}0\,\% & non \\
\texttt{A15\_D150\_C00} & 15\,\% & $\times$1{,}50 & -- & 1\,324\,625 & 188 & 1\,606 & 37{,}8\,\% & respectées \\
\texttt{A15\_D150\_C05} & 15\,\% & $\times$1{,}50 & $-$5\,\% & 1\,362\,545 & 227 & 1\,468 & 40{,}0\,\% & respectées \\
\texttt{A15\_D150\_C10} & 15\,\% & $\times$1{,}50 & $-$10\,\% & 1\,394\,585 & 262 & 1\,342 & 42{,}0\,\% & respectées \\
\texttt{A15\_D150\_C20} & 15\,\% & $\times$1{,}50 & $-$20\,\% & 1\,479\,206 & 329 & 1\,139 & 45{,}9\,\% & non \\
\texttt{A15\_D200\_C00} & 15\,\% & $\times$2{,}00 & -- & 1\,984\,825 & 434 & 1\,624 & 37{,}6\,\% & respectées \\
\texttt{A15\_D200\_C05} & 15\,\% & $\times$2{,}00 & $-$5\,\% & 2\,026\,705 & 473 & 1\,486 & 39{,}7\,\% & respectées \\
\texttt{A15\_D200\_C10} & 15\,\% & $\times$2{,}00 & $-$10\,\% & 2\,060\,215 & 508 & 1\,357 & 41{,}7\,\% & respectées \\
\texttt{A15\_D200\_C20} & 15\,\% & $\times$2{,}00 & $-$20\,\% & 2\,152\,276 & 575 & 1\,157 & 45{,}5\,\% & non \\
\texttt{A20\_D100\_C00} & 20\,\% & $\times$1{,}00 & -- & 733\,875 & 4 & 1\,395 & 39{,}9\,\% & respectées \\
\texttt{A20\_D100\_C05} & 20\,\% & $\times$1{,}00 & $-$5\,\% & 747\,915 & 14 & 1\,332 & 41{,}0\,\% & respectées \\
\texttt{A20\_D100\_C10} & 20\,\% & $\times$1{,}00 & $-$10\,\% & 771\,465 & 26 & 1\,290 & 41{,}8\,\% & respectées \\
\texttt{A20\_D100\_C20} & 20\,\% & $\times$1{,}00 & $-$20\,\% & 821\,455 & 79 & 1\,091 & 45{,}9\,\% & non \\
\texttt{A20\_D125\_C00} & 20\,\% & $\times$1{,}25 & -- & 1\,025\,824 & 73 & 1\,611 & 36{,}6\,\% & respectées \\
\texttt{A20\_D125\_C05} & 20\,\% & $\times$1{,}25 & $-$5\,\% & 1\,069\,384 & 105 & 1\,512 & 38{,}1\,\% & respectées \\
\texttt{A20\_D125\_C10} & 20\,\% & $\times$1{,}25 & $-$10\,\% & 1\,095\,464 & 139 & 1\,382 & 40{,}2\,\% & respectées \\
\texttt{A20\_D125\_C20} & 20\,\% & $\times$1{,}25 & $-$20\,\% & 1\,176\,115 & 206 & 1\,177 & 44{,}0\,\% & non \\
\texttt{A20\_D150\_C00} & 20\,\% & $\times$1{,}50 & -- & 1\,348\,124 & 188 & 1\,653 & 36{,}0\,\% & respectées \\
\texttt{A20\_D150\_C05} & 20\,\% & $\times$1{,}50 & $-$5\,\% & 1\,386\,044 & 227 & 1\,515 & 38{,}0\,\% & respectées \\
\texttt{A20\_D150\_C10} & 20\,\% & $\times$1{,}50 & $-$10\,\% & 1\,416\,084 & 262 & 1\,385 & 40{,}1\,\% & respectées \\
\texttt{A20\_D150\_C20} & 20\,\% & $\times$1{,}50 & $-$20\,\% & 1\,499\,705 & 329 & 1\,180 & 44{,}0\,\% & non \\
\texttt{A20\_D200\_C00} & 20\,\% & $\times$2{,}00 & -- & 2\,008\,823 & 434 & 1\,672 & 35{,}8\,\% & respectées \\
\texttt{A20\_D200\_C05} & 20\,\% & $\times$2{,}00 & $-$5\,\% & 2\,050\,703 & 473 & 1\,534 & 37{,}8\,\% & respectées \\
\texttt{A20\_D200\_C10} & 20\,\% & $\times$2{,}00 & $-$10\,\% & 2\,081\,714 & 508 & 1\,400 & 39{,}9\,\% & respectées \\
\texttt{A20\_D200\_C20} & 20\,\% & $\times$2{,}00 & $-$20\,\% & 2\,172\,775 & 575 & 1\,198 & 43{,}6\,\% & non \\
\end{longtable}
}

